% Generated draft: requires mathematical and source review.
\documentclass[11pt]{article}
\usepackage[margin=25mm]{geometry}
\usepackage{fontspec}
\setmainfont{Latin Modern Roman}
\usepackage{amsmath,amssymb,mathtools,unicode-math}
\setmathfont{Latin Modern Math}
\usepackage{enumitem,needspace,xurl,hyperref}
\setlength{\parindent}{0pt}
\setlength{\parskip}{5pt}
\setlength{\emergencystretch}{3em}
\newcommand{\N}{\mathbb{N}}
\newcommand{\Z}{\mathbb{Z}}
\newcommand{\Q}{\mathbb{Q}}
\newcommand{\R}{\mathbb{R}}
\newcommand{\C}{\mathbb{C}}
\newcommand{\F}{\mathbb{F}}
\newcommand{\A}{\mathbb{A}}
\newcommand{\PP}{\mathbb{P}}
\newcommand{\E}{\mathbb{E}}
\newcommand{\eps}{\varepsilon}
\newcommand{\dd}{\,\mathrm{d}}
\begin{document}
\section*{UnsolvedMath ID: 11000326}
\textbf{Source problem code:} catalogue entry
\subsection{Definitions and mathematical statement}
For a surface $\Sigma$ in the source's convention, a homology cylinder is represented by an oriented cobordism $C$ with two marked boundary copies $i_\pm:\Sigma\to C$ whose inclusions induce homology isomorphisms, with the appropriate agreement along boundary. Conventions differ on whether $i_{+*}=i_{-*}$ is required; the excerpt does not supply its definition. The usual stacking operation first gives a monoid. Quotienting by the intended homology-cobordism relation gives the group mentioned in the question, which is not the unquotiented cylinder monoid.

Homology cobordism conventionally requires a marked $4$-dimensional cobordism between the cylinders inducing homology isomorphisms and respecting their boundary data. The smooth versus topological category, coefficient ring, surface boundary data and exact equivalence relation are not specified here.

Problem 11.5 asks to generalize Hain's infinitesimal presentation of the Torelli Lie algebra [29] to this homology-cobordism group. An answer must identify the corresponding infinitesimal Lie algebra and describe its generators and relations in the intended analogy. The excerpt does not provide Hain's presentation, its genus hypotheses, the filtration, coefficient field or completion. A rational Malcev Lie algebra, an associated lower-central-series Lie algebra and a Johnson-filtration Lie algebra need not coincide. No finite presentation of the cylinder group or of its Lie algebra is asserted or imposed merely by the word generalize.
\subsection{Short English statement}
Extend Hain's infinitesimal Torelli presentation to the group obtained from homology cylinders by homology cobordism, after fixing the cylinder, category and Lie-algebra conventions.
\subsection{Sources}
\begin{itemize}
\item Local registry: \url{data/unsolvedmath/problems.json}, permanent UnsolvedMath ID $11000326$, exact source code \texttt{catalogue entry}.
\item Farb, \emph{Problems on Mapping Class Groups and Related Topics}, Problem 11.5, PDF page 379, chapter ``Cohomological structure of the mapping class group and beyond'': \url{https://www.math.uchicago.edu/~farb/papers/mcgbook.pdf}. Supplied access date: 2026-07-29; extraction: pdf-font-structure; page-disambiguated numbering among 329 headings.
\item Hain [29] for the infinitesimal Torelli presentation; full bibliography entry, theorem and preceding homology-cylinder definitions are absent. The registry reports no such generalization known to its auditor, without an additional primary citation.
\end{itemize}
\end{document}
